% 原书物理页83--87，2.2.
\section{代数闭域与代数闭包}\label{sec:galois-2-2}
\subsection{定义}\label{sec:galois-2-2-1}
\begin{definition*}
若域 $K$ 满足下列等价条件，则称 $K$ 为\newterm{galois-algebraically-closed}{代数闭域}{algebraically closed field}：
\begin{itemize}
\item 每个多项式 $f\in K[X]$ 在 $K$ 中都有一个根；
\item 每个多项式 $f\in K[X]$ 都分裂；
\item $K[X]$ 的不可约元素恰为一次多项式.
\end{itemize}
\end{definition*}
\noindent\textbf{译注.} 原文前两项未排除非零常数多项式；按通常约定，应对非常数多项式理解这些条件.

回忆，``$f$ 在 $K$ 中的一个根''是满足 $(X-x)\mid f$ 的元素 $x\in K$（等价地，$f(x)=0$）；``$f$ 分裂''是指 $f$ 可分解为 $f=\lambda(X-\alpha_1)\cdots(X-\alpha_n)$. 也回忆下面的著名定理.
\begin{theorem*}
$\mathbb C$ 是代数闭域.
\end{theorem*}
\begin{proof}
设 $f=X^n+a_1X^{n-1}+\cdots+a_{n-1}X+a_n$. 反证，假设 $f$ 在 $\mathbb C$ 上不取零值. 我们证明 $f$ 在 $\mathbb C$ 上达到其最小值. 有 $f(0)=a_n\ne0$. 由于 $\lim_{|z|\to+\infty}|P(z)|=+\infty$，存在 $R>0$，使得 $|z|>R\Rightarrow |P(z)|>|a_n|$. 圆盘 $D=\{z\in\mathbb C,\ |z|\le R\}$ 紧，故 $f$ 在其上达到最小值. 这个最小值不大于 $f(0)=|a_n|$，因而也是 $f$ 在 $\mathbb C$ 上的最小值.

必要时作变量代换 $X\mapsto X+z_0$，可设 $f$ 在 $0$ 处达到最小值. 再乘以适当常数，可设 $f(0)=1$. 若 $k$ 是 $f-1$ 在 $0$ 处的消失阶，则在 $z=0$ 附近有
\begin{align*}
f(z)=1+a_{n-k}z^k+o(z^k).
\end{align*}
取 $\theta$ 使 $a_{n-k}e^{ik\theta}=-|a_{n-k}|$（即 $k\theta=\pi-\arg(a_{n-k})$）. 则 $f(re^{i\theta})=1-|a_{n-k}|r^k+o(r^k)$，从而当 $r$ 非零且充分小时，$f(re^{i\theta})<1$，与 $1$ 为 $f$ 的最小值矛盾.
\end{proof}
\noindent\textbf{译注.} 此证明依原文保留记号：$P$ 未定义，应对应 $f$；文中的``$f$ 的最小值''以及末尾不等式，应就模 $|f|$ 理解，复数本身无此大小关系. 原文亦写作 $f(0)=|a_n|$，未区分函数值与模.

\subsection{定义}\label{sec:galois-2-2-2}
\begin{definition*}
设 $k$ 是域. $k$ 的一个（绝对）\newterm{galois-algebraic-closure}{代数闭包}{algebraic closure}是一个代数扩张 $k\subset\bar k$，其中 $\bar k$ 为代数闭域.
\end{definition*}
\begin{lemma*}
设 $k\subset K$ 是域扩张，且 $K$ 为代数闭域. 则 $k$ 在 $K$ 中的（相对）代数闭包 $K_{\mathrm{alg}}$ 是 $k$ 的一个（绝对）代数闭包.
\end{lemma*}
\begin{proof}
由构造，$K_{\mathrm{alg}}$ 在 $k$ 上代数. 故只须证明 $K_{\mathrm{alg}}$ 代数闭. 设 $f\in K_{\mathrm{alg}}[X]$. 因为 $K$ 代数闭，$f$ 在 $K$ 中有根 $\alpha$. 这个元素在 $K_{\mathrm{alg}}$ 上代数，因而也在 $k$ 上代数. 所以它属于 $K_{\mathrm{alg}}$.
\end{proof}
\begin{example*}
$\overline{\mathbb Q}$ 是 $\mathbb Q$ 的一个代数闭包.
\end{example*}
给定两个扩张 $k\subset K$ 与 $k\subset K'$，一个扩张态射是一个 $k$ 线性的域态射 $K\to K'$. 因而它在 $k$ 上诱导恒等映射. 这样的态射总是单射，故也称为嵌入.

\subsection{命题}\label{sec:galois-2-2-3}
\begin{proposition*}
若 $k\subset\bar k$ 是 $k$ 的一个代数闭包，则 $k$ 的任意代数扩张均可嵌入 $\bar k$.
\end{proposition*}
\begin{proof}
设 $K$ 为 $k$ 的代数扩张，$K'\subset K$ 是一个配有嵌入 $\iota:K'\hookrightarrow\bar k$ 的子扩张. 关键在于，对任意 $\alpha\in K$，嵌入 $\iota$ 都能延拓到 $K'(\alpha)$. 事实上，$\alpha$ 在 $k$ 上代数，当然也在 $K'$ 上代数，故
\begin{align*}
K'(\alpha)=K'[\alpha]\simeq K'[X]/(f_\alpha),
\end{align*}
其中 $f_\alpha\in K'[X]$ 表示 $\alpha$ 在 $K'$ 上的极小多项式. 对 $f_\alpha$ 的系数施以 $\iota$，得到多项式 $\iota(f_\alpha)\in\bar k[X]$. 它是 $f_\alpha$ 在唯一的 $K'$ 代数态射 $K'[X]\to\bar k[X]$ 下的像；这个态射将 $X$ 映到 $X$，在 $K'$ 上则由 $\iota$ 映入 $\bar k$. 由于 $\bar k$ 代数闭，可以在 $\bar k$ 中选择 $\iota(f_\alpha)$ 的一个根 $x$. 考虑将 $X$ 映到 $x$ 且延拓 $\iota$ 的唯一 $K'$ 代数态射 $\varphi:K'[X]\to\bar k$. 对任意 $f\in K'[X]$，有 $\varphi(f)=\iota(f)(x)$. 特别地，$\varphi(f_\alpha)=0$，故 $\varphi$ 通过一个 $K'$ 代数态射
\begin{align*}
K'[X]/(f_\alpha)\longrightarrow\bar k
\end{align*}
分解. 这个态射必为域嵌入，并如所要求地延拓 $\iota$. 但要注意，这个延拓远非典范的，因为它依赖于所选的 $f_\alpha$ 的根 $x$，甚至依赖于 $\alpha$，因为 $K'(\alpha)$ 当然还有其他生成元.

我们借助 Zorn 引理来处理延拓不唯一的问题. 考虑所有二元组 $(K',\iota')$ 组成的集合 $\mathcal P$，其中 $K'\subset K$ 是 $k$ 的子扩张，$\iota':K'\hookrightarrow\bar k$ 是嵌入. 以如下关系赋予偏序：
\begin{align*}
(K',\iota')\le(K'',\iota'')\quad\Longleftrightarrow\quad K'\subset K''\text{ 且 }\iota'=\iota''|_{K'}.
\end{align*}
此序是``归纳的''，即每个递增序列都有上界. 事实上，若 $(K'_n,\iota'_n)_{n\in\mathbb N}$ 是递增序列，则 $K':=\bigcup_nK'_n$ 为子域；令 $\iota'(x)=\iota'_n(x)$，此值与满足 $x\in K'_n$ 的 $n$ 的选取无关，从而定义了嵌入 $\iota'$. 二元组 $(K',\iota')$ 就是所有 $(K'_n,\iota'_n)$ 的上界. Zorn 引理断言，每个归纳有序集都具有极大元. 取 $\mathcal P$ 的极大元 $(K',\iota')$. 若存在 $\alpha\in K\setminus K'$，证明开头的构造便与其极大性矛盾. 所以 $K'=K$.
\end{proof}
\noindent\textbf{译注.} 原文只验证了可数递增序列有上界；应用 Zorn 引理所需的条件是每条链都有上界. 另，前段所谓 $f_\alpha$ 的根 $x$，在目标域中准确地说是 $\iota(f_\alpha)$ 的根.
\begin{corollary*}
$k$ 的任意两个代数闭包 $\bar k$ 与 $\bar k'$ 作为 $k$ 的扩张是同构的.
\end{corollary*}
\begin{proof}
由命题，存在嵌入 $\bar k'\hookrightarrow\bar k$. 此嵌入的像 $K$ 同构于 $\bar k'$，因而代数闭. $\bar k$ 的每个元素 $\alpha$ 都在 $k$ 上代数，当然也在 $K$ 上代数. 因 $K$ 代数闭，$\alpha$ 在 $K$ 上的极小多项式 $f_\alpha$ 为一次，故形如 $X-a_0$. 于是 $\alpha=a_0\in K$，进而 $K=\bar k$ 且 $\bar k'\xrightarrow{\sim}\bar k$.
\end{proof}
\begin{remark*}
通常不存在典范同构. 例如 $\mathbb C$、$\mathbb R[X]/(X^2+1)$ 和 $\mathbb R[X]/(X^2+X+1)$ 都是 $\mathbb R$ 的代数闭包，但这些域之间没有典范同构. 推论也可以表述为：代数闭包在不唯一的同构意义下唯一.
\end{remark*}
\begin{remark*}
设 $K$ 为 $k$ 的有限扩张，假定 $K$ 是单生成的，取 $\alpha\in K$ 使 $K=k(\alpha)$. 记 $f_\alpha\in k[X]$ 为 $\alpha$ 的极小多项式，则命题的证明给出双射 $\iota\mapsto\iota(\alpha)$：
\begin{align*}
\{k\text{ 线性嵌入 }\iota:K\hookrightarrow\bar k\}\longleftrightarrow\{f_\alpha\text{ 在 }\bar k\text{ 中的根}\}.
\end{align*}
\end{remark*}

\subsection{代数闭包的构造}\label{sec:galois-2-2-4}
现在证明每个域 $k$ 都有代数闭包. 先介绍一种归纳构造域的方法.
\begin{lemma*}
设 $k_0\xrightarrow{\tau_0}k_1\xrightarrow{\tau_1}\cdots\xrightarrow{\tau_{n-1}}k_n\xrightarrow{\tau_n}\cdots$ 为一列域态射. 则存在域 $k_\infty$ 及嵌入 $\iota_n:k_n\hookrightarrow k_\infty$，使每个 $n>0$ 均有 $\iota_n\circ\tau_{n-1}=\iota_{n-1}$，并满足以下泛性质：对每个域 $K$ 以及满足 $\sigma_n\circ\tau_{n-1}=\sigma_{n-1}$（$n>0$）的任意一族嵌入 $\sigma_n:k_n\to K$，存在唯一嵌入 $\sigma:k_\infty\to K$，使得对每个 $n$ 均有 $\sigma\circ\iota_n=\sigma_n$.
\end{lemma*}
\begin{proof}
若给定的态射列是某个``大''域 $K$ 内的一列包含 $k_0\subset k_1\subset\cdots\subset k_n\subset\cdots$，只须取 $k_\infty:=\bigcup_nk_n$. 微妙之处在于，这里给定的是未包含在一个大域中的``抽象''域，因而必须从外部构造``它们的并''.

考虑直和 $\bigoplus_{n\in\mathbb N}k_n$. 这是一个 $k_0$ 向量空间，甚至是一个无单位元的 $k_0$ 代数（它是 $\prod_{n\in\mathbb N}k_n$ 的一个理想）. 按直和定义，有嵌入 $\tilde\iota_n:k_n\hookrightarrow\bigoplus_{m\in\mathbb N}k_m$，将 $x_n\in k_n$ 映到仅第 $n$ 项为 $x_n$、其余项均为零的序列. 令 $R$ 为下列元素生成的向量子空间：
\begin{align*}
\tilde\iota_n(\tau_{n-1}(x_{n-1}))-\tilde\iota_{n-1}(x_{n-1})
=(0,\ldots,0,-x_{n-1},\tau_{n-1}(x_{n-1}),0,0,\ldots),
\end{align*}
其中 $n\in\mathbb N^*$，$x_{n-1}\in k_{n-1}$. 置
\begin{align*}
k_\infty:=\left(\bigoplus_n k_n\right)/R,\qquad
\iota_n:k_n\xrightarrow{\tilde\iota_n}\bigoplus_{m\in\mathbb N}k_m\twoheadrightarrow k_\infty.
\end{align*}
由定义，每个 $n>0$ 都满足 $\tilde\iota_{n-1}(k_{n-1})\subset\tilde\iota_n(k_n)+R$，故 $\iota_{n-1}(k_{n-1})\subset\iota_n(k_n)$. 又 $k_\infty=\sum_{n\in\mathbb N}\iota_n(k_n)$，因此 $k_\infty=\bigcup_n\iota_n(k_n)$. 另一方面，由各 $\tau_i$ 的单射性可知，$R$ 中每个序列 $(x_m)_{m\in\mathbb N}$ 至少有两项非零. 从而对每个 $n$ 都有 $\tilde\iota_n(k_n)\cap R=\{0\}$，故 $\iota_n$ 单射. 因此 $\iota_n$ 给出 $k_n$ 到其像 $\iota_n(k_n)$ 的 $k_0$ 线性同构，且在这些同构下，每个包含 $\iota_{n-1}(k_{n-1})\subset\iota_n(k_n)$ 都对应于态射 $\tau_{n-1}$. 换言之，下图交换：
\[
\begin{tikzcd}
k_n\arrow[r,"\sim" above,"\iota_n" below]&\iota_n(k_n)\\
k_{n-1}\arrow[u,"\tau_{n-1}"]\arrow[r,"\sim" above,"\iota_{n-1}" below]&\iota_{n-1}(k_{n-1})\arrow[u,hook]
\end{tikzcd}
\]
现在可给 $k_0$ 向量空间 $k_\infty$ 定义乘法：对 $x,y\in k_\infty$，取充分大的 $n$ 使 $x,y\in\iota_n(k_n)$，置 $xy:=\iota_n(\iota_n^{-1}(x)\iota_n^{-1}(y))$. 此定义与 $n$ 的选取无关，并使 $k_\infty$ 成为 $k_0$ 代数. 因为此代数是子域 $\iota_n(k_n)$ 的并，它本身也是域.

$k_\infty$ 连同 $\iota_n$ 的泛性质也以同样的方式证明：若 $x\in k_\infty$ 属于 $\iota_n$ 的像，就置 $\sigma(x):=\sigma_n(\iota_n^{-1}(x))$；整个构造保证了这与 $n$ 的选取无关. $\sigma$ 的唯一性来自 $k_\infty=\bigcup_n\iota_n(k_n)$.
\end{proof}
\noindent\textbf{译注.} 原文称 $R$ 中``每个序列''至少有两个非零项，须排除零序列；此处保留原表述.

我们将在下一个引理的假设下应用上述引理.
\begin{lemma*}
沿用前一引理的记号. 假设对每个 $n\ge0$ 及每个多项式 $f_n\in k_n[X]$，多项式 $\tau_n(f_n)\in k_{n+1}[X]$ 在 $k_{n+1}$ 中都有一个根. 则 $k_\infty$ 是代数闭域.
\end{lemma*}
\begin{proof}
设 $f\in k_\infty[X]$. 存在 $n$，使 $f$ 的系数都在 $\iota_n(k_n)$ 中. 于是存在（唯一的）多项式 $f_n\in k_n[X]$，使 $f=\iota_n(f_n)$. 由假设，多项式 $\tau_n(f_n)\in k_{n+1}[X]$ 在 $k_{n+1}$ 中有根 $x_{n+1}$. 因此 $\iota_{n+1}(x_{n+1})$ 是多项式 $\iota_{n+1}(\tau_n(f_n))=\iota_n(f_n)=f$ 在 $k_\infty$ 中的根. 故 $k_\infty$ 代数闭.
\end{proof}
因此，为证明代数闭包的存在性，只须归纳应用下面的命题.
\begin{proposition*}
设 $k$ 为域. 存在 $k$ 的一个代数扩张 $K$，使每个多项式 $f\in k[X]$ 在其中都有一个根.
\end{proposition*}
\begin{remark*}[{根域}]
在给出证明前，先指出：对一个给定的不可约多项式 $f\in k[X]$，容易构造 $k$ 的扩张 $K_f$，使 $f$ 在其中有根. 只须取 $K_f:=k[X]/(f)$；因 $(f)$ 为极大理想，这是一个域，且 $X$（更确切地说，它的像）是 $f$ 的根. 这样的域 $K_f$ 称为 $f$ 的\newterm{galois-root-field}{根域}{root field}.
\end{remark*}
同样容易归纳构造 $k$ 的扩张 $K_{f_1,\ldots,f_n}$，使每个给定的多项式 $f_i$ 在其中都有根. 对一个多项式族 $(f_n)_{n\in\mathbb N}$，也可如此处理，例如使用第一个引理. 但一般而言，不可约多项式的集合未必可数. 下面的证明将这一想法推广到一般情形.
\begin{proof}[命题的证明]
记 $(f_i)_{i\in I}$ 为 $k[X]$ 中所有首一不可约多项式构成的族. 考虑以 $I$ 为不定元指标集的多项式环 $\mathcal R:=k[(X_i)_{i\in I}]$，以及其中由 $f_i(X_i)$（$i\in I$）生成的理想 $\mathcal I$.

假设这个理想是真理想. 则由 Zorn 引理，它包含于 $\mathcal R$ 的某个极大理想 $\mathfrak m$. 商 $K:=\mathcal R/\mathfrak m$ 是包含 $k$ 的域. 按构造，$X_i$ 在 $K$ 中的像是 $f_i$ 的根. 又 $K$ 作为扩张由 $X_i$ 的像生成，故 $K$ 是代数扩张并满足命题的要求.

因此只须证明 $\mathcal I$ 确为 $\mathcal R$ 的真理想. 反证，假设 $\mathcal I=\mathcal R$. 则存在有限子集 $J\subset I$ 及元素 $g_j\in\mathcal R$，使得 $\sum_{j\in J}g_jf_j(X_j)=1$. 因 $J$ 有限，由前面的讨论，存在 $k$ 的扩张 $K_J$，使每个 $f_j$ 都有根，记作 $x_j$. 考虑唯一的 $k$ 代数态射 $\mathcal R\to K_J$，当 $i\in J$ 时将 $X_i$ 映到 $x_i$，否则映到 $0$. 它将 $f_j(X_j)$ 映到 $f_j(x_j)=0$，因而也将 $\sum_{j\in J}g_jf_j(X_j)$ 映到 $0$. 由于域 $K_J$ 中 $0\ne1$，得到矛盾.
\end{proof}
\noindent\textbf{译注.} 本分目中第二引理及命题的``每个多项式''均沿用原文；非零常数多项式须排除，与本节开头相同.

\subsection{定理}\label{sec:galois-2-2-5}
\begin{theorem*}
每个域都有代数闭包.
\end{theorem*}
\begin{proof}
设 $k=k_0$ 为域. 前一命题给出代数扩张 $k_1$，使每个 $f\in k_0[X]$ 在其中有根. 归纳可得一列代数扩张 $k_0\subset k_1\subset\cdots\subset k_n\subset\cdots$，满足上面第二个引理的假设. 第一个引理的构造便给出包含 $k$ 且在 $k$ 上代数的代数闭域 $k_\infty$.
\end{proof}



